\documentclass[12pt]{article}
\renewcommand{\baselinestretch}{1.5}

%%%%%%%%%%packages
\usepackage[english]{babel}
%%%%%%%%%%%%%%%%%%%%%%%%%%%%%%%%%%%%%%%%%
%%%%%%%%%%%%%%%%%%%%%%%%%%%%%%%%%%%%%%%%
\usepackage[latin1]{inputenc}
%%%%%%%%%%%%%%%%%%%%%%%%%%%%%%%%%%%%%%%%
%%%%%%%%%%%%%%%%%%%%%%%%%%%%%%%%%%%%%%%%%
\PassOptionsToPackage{usenames,dvipsnames}{xcolor}
\usepackage[dvips]{graphicx,psfrag,overpic,color} %Necesario para figuras
\usepackage{latexsym} %Necesario para simbolos ampliados
\usepackage{amssymb} %Necesario para simbolos AMS
\usepackage{amsmath}%Necesario para simbolos AMS extra
\usepackage{amsfonts}%Necesario para fuentes AMS extra
\usepackage{amsthm}
\usepackage{pstricks}  % since the dash is rendered by pstricks!
\usepackage[postscript]{ucs}
\usepackage{fancyhdr}
\usepackage{listings}
\usepackage[a4paper,left=3cm,right=3cm,top=2.5cm,bottom=2.5cm]{geometry}
\usepackage{url}
\usepackage{pdfpages}
\usepackage{mathtools}
\usepackage{titling}
\usepackage[nottoc,numbib]{tocbibind}
\usepackage{mathptmx}
\usepackage[explicit]{titlesec}


\titleformat{\section}[block]{\filcenter}{\thesection. \MakeUppercase{#1}}{1em}{}
\titleformat{name=\section, numberless}[block]{\filcenter}{\MakeUppercase{#1}}{1em}{}
\titleformat{\subsection}[hang]{}{}{1em}{\thesubsection. #1}
\titleformat{\subsubsection}[hang]{}{}{1em}{{\normalfont\thesubsubsection. \textit{#1}}}


\newtheorem{theorem}{{Theorem}}
\newtheorem{prop}{{Proposition}}
\newtheorem{defi}{{Definition}}
\newtheorem{claim}{{Claim}}

\makeindex

\pretitle{%
  \begin{center}
  \LARGE
  \vspace{1in}
  \includegraphics[width=4.56in]{logo_uab2}\\[\bigskipamount]
  \vspace{1.5in}
}
\posttitle{\end{center}}

\begin{document}


\selectlanguage{english}

\title{{\large TITLE: Merger Simulations: A Theoretical and Practical Approach\\
AUTHOR: Andrea Dana Klein Villalba\\
DEGREE: Economics\\
ADVISOR: Jordi Perdiguero Garc\'ia\\
DATE: June 7th, 2019\\}}
\date{}

%-------------------------------------------------------------%
\clearpage\maketitle
\thispagestyle{empty}
\newpage

\begin{abstract}
Hello, it's me,
\end{abstract}
\newpage

%--------------------------------------------------------------
% Taula de continguts
% S'elabora automaticament
%-------------------------------------------------------------
 \tableofcontents
 \newpage
%------------------------------------------------------------
% Cos del treball
%Es poden fer servir Parts, Capitols, Seccions i Subseccions
%-------------------------------------------------------------
\pagenumbering{arabic}

\section{Introduction}



\subsection{Patata}
sdfsfs


\section{Models}

The models used in this project are characterized by Bertrand Competition, where the suppliers choose prices according to demand. This will be modeled as an economy in which there are $N$ products and $M$ different firms that produce them.

\begin{defi}
The set of products is denoted by $\mathcal{N} = \{1, 2, \dots, N\}$ and the subset of products produced by the same firm will be written as $\mathcal{N}_m \subseteq \mathcal{N}$. The set of firms will be written as $\mathcal{M}$.
\end{defi}

It is important to note that each product is exactly produced by one firm, i.e.  $\mathcal{N}_{m_1} \cap \mathcal{N}_{m_2} = \emptyset$ if $m_1$ and $m_2$ are two different firms and $\bigcup\limits_{m\in\mathcal{M}} \mathcal{N}_m = \mathcal{N}$. The firms in this economy will be profit-maximizing firms, and therefore we can write their profit function as
\begin{equation}
\pi_m(p) = \sum_{j \in \mathcal{N}_m} (p_j - c_j) \cdot q_j(p) \label{profit_eq}
\end{equation}

Where $p = (p_1, p_2, \dots, p_N)^t$ is the column vector of prices, and $c_j = \frac{\partial C_j}{\partial q_j}$ are the marginal costs and $q_j$ is the demand for the $j$-th good. 

\begin{defi}
The demand function $q$ is
\begin{eqnarray*}
q: ~ \mathbb{R}^N & \rightarrow & \mathbb{R}^N \\
p & \mapsto & q(p)
\end{eqnarray*}
where $q(p) = (q_1(p), q_2(p), \dots, q_N(p))^t$ is the column vector of the demands.
\end{defi}

\begin{defi}
The marginal cost function $c$ is
\begin{eqnarray*}
c: ~ \mathbb{R}^N & \rightarrow & \mathbb{R}^N \\
q & \mapsto & c(q)
\end{eqnarray*}
where $c(q) = (c_1(q_1), c_2(q_2), \dots, c_N(q_N))^t$ is the column vector of the marginal costs.
\end{defi}

Notice that, while $c(q)$ is a function of $q$, the latter is also a function of $p$, so $c(q)$ is a function of $p$ as well. 

For the given profit function, it is maximized when the first order conditions are fulfilled. Such that, deriving from equation \ref{profit_eq} with respect to $p_k$ such that $k \in \mathcal{N}_m$.
\begin{eqnarray}
\frac{\partial \pi_m}{\partial p_k} &=& \frac{\partial}{\partial p_k} \sum_{j \in \mathcal{N}_m} (p_j - c_j(q_j)) \cdot q_j(p) \\
&=& \sum_{j \in \mathcal{N}_m} \frac{\partial}{\partial p_k} \left[(p_j - c_j(q_j)) \cdot q_j(p)\right] \\
&=& \sum_{j \in \mathcal{N}_m} q_j(p) \frac{\partial (p_j - c_j)}{\partial p_k}(q_j) + \sum_{j \in \mathcal{N}_m} \frac{\partial q_j}{\partial p_k}(p) \cdot (p_j - c_j(q_j))\\
&=& q_k(p) - \sum_{j \in \mathcal{N}_m} q_j(p) \frac{\partial c_j}{\partial p_k} (q_j) + \sum_{j \in \mathcal{N}_m} \frac{\partial q_j}{\partial p_k}(p) \cdot (p_j - c_j(q_j)) = 0 \label{profit_FOC_eq}
\end{eqnarray}

Since each product is produced by exactly one firm, deriving as so for each product, we obtain a system of $N$ equations that is solved for the optimal prices $p^*$. In order to solve this system, it is useful to introduce the following variables:
\begin{defi}
Let $D = (d_{ij})$ be the ownership matrix of size $N \times N$, such that\\
\begin{center}
$d_{ij} = \left\{
\begin{tabular}{ll}
1 & if $i,j$ produced by the same firm\\
0 & otherwise
\end{tabular}\right.$
\end{center}
\end{defi}

By using the ownership matrix $D$, we can rewrite the first order condition from equation \ref{profit_FOC_eq} as follows
\begin{equation}
q_k(p) - \sum_{j=1}^N q_j(p) \frac{\partial c_j}{\partial p_k} (q_j)d_{jk} + \sum_{j=1}^N \frac{\partial q_j}{\partial p_k}(p) \cdot (p_j - c_j(q_j))d_{jk} = 0 \label{profit_FOC_with_Ds_eq}
\end{equation}


\subsection{Linear Demands with Constant Marginal Costs}

In this model we will assume that demands are linear, such that

\begin{equation*}
q_k(p) = a_k + \sum_{j=1}^N b_{kj}p_j
\end{equation*}

Because marginal costs are constant, the costs function $C_j$ will be characterized by
\begin{equation*}
C_j (q_j)= f_j + v_j \cdot q_j
\end{equation*}

Where $f_j$ is measure of fixed costs, $v_j$ is a measure of variable costs, and  $j \in \mathcal{N}$. Thus, marginal costs are
\begin{equation*}
c_j = \frac{\partial C_j}{\partial q_j}(q_j) = v_j
\end{equation*}

Computing the first order conditions from equation \ref{profit_FOC_with_Ds_eq}, we obtain
\begin{eqnarray*}
\frac{\partial\pi_m(p)}{\partial p_k}
&=& q_k(p) + \sum_{j = 1}^N d_{jk}b_{jk} \cdot (p_j - c_j) = 0
\end{eqnarray*}

Because we have one such first order condition for each product, we can write the system in a more compact way using matrices and vectors.
\begin{equation*}
a + B^t p + (D \circ B)(p - c) = 0
\end{equation*}

Where $\circ$ denotes the Hadamard product.
Isolating $p$, we obtain the optimum prices $p^*$
\begin{equation*}
p^* = (B^t + (D \circ B))^{-1}((D \circ B) c - a)
\end{equation*}

Then, optimum quantities are defined by
\begin{equation*}
q^* = a + B^t \cdot p^*
\end{equation*}

We are now interested in computing consumer welfare $CS_j$ for each product. 
\begin{equation*}
CS_j = \int_{p_j^*}^{p_j(-\frac{a_j}{b_{jk})}} p_j d p_j = \frac{1}{2} \cdot [(p_j(-\frac{a_j}{b_{jk}}) - p_j^*) \cdot q_j^*]
\end{equation*}

\subsection{Linear Demands with Linear Marginal Costs}

Let $N$ be the number of products, and $\mathcal{N} = \{1, 2, \dots, N\}$.

Let $q_k (p)$ be the demand of the $k$-th good. Because it is a linear demand, its curve is defined as

\begin{equation*}
q_k(p) = a_k + \sum_{j=1}^N b_{kj}p_j
\end{equation*}

Let $m$ be a firm that produces $\mathcal{N}_m \subseteq \mathcal{N}$. Then, its profit is defined by

\begin{equation*}
\pi_m(p) = \sum_{j \in \mathcal{N}_m} (p_j - c_j)\cdot q_j(p)
\end{equation*}

The costs $C_j$ are characterized by

\begin{equation*}
C(q_j) = f + v_j \cdot q_j + \frac{w_j}{2} \cdot q_j ^ {2}
\end{equation*}

Where $j \subseteq \mathcal{N}_m$, $f$ is a constant fixed cost of production, $v$ is the variable cost and $w$ is a quadratic variable cost arising from the production of an additional unit of $q_j$. Thus, marginal costs are

\begin{equation*}
c_j = \frac{\partial C(q_j)}{\partial q_j}= v_j + w_j \cdot q_j.
\end{equation*}
The firm, then, seeks to maximize its profit by setting their prices according to the following formula

\begin{equation*}
\max_{\{p_j ~ | ~ j \in \mathcal{N}_m\}} \pi_m(p)
\end{equation*}

If we derive with respect to $p_k$, where $k \subseteq \mathcal{N}_m$, then we obtain

\begin{eqnarray*}
\frac{\partial\pi_m(p)}{\partial p_k} &=& \frac{\partial}{\partial p_k} \sum_{j \in \mathcal{N}_m} (p_j - c_j) \cdot q_j(p) \\
&=& \sum_{j \in \mathcal{N}_m} \frac{\partial}{\partial p_k} \left[(p_j - c_j) \cdot q_j(p)\right] \\
&=& \sum_{j \in \mathcal{N}_m} \frac{\partial (p_j - c_j)}{\partial p_k} \cdot q_j(p) + \sum_{j \in \mathcal{N}_m} \frac{\partial q_j(p)}{\partial p_k} \cdot (p_j - c_j)\\
&=& q_k(p) + \sum_{j \in \mathcal{N}_m} b_{jk} \cdot p_j - \sum_{j \in \mathcal{N}_m} b_{jk} \cdot v_j - 2 \cdot \sum_{j \in \mathcal{N}_m} b_{jk} \cdot w_j \cdot q_j (p) 
\end{eqnarray*}

From here we obtain the first order condition

\begin{equation*}
q_k(p) + \sum_{j \in \mathcal{N}_m} b_{jk} \cdot p_j - \sum_{j \in \mathcal{N}_m} v_j \cdot b_{jk} - 2 \cdot \sum_{j \in \mathcal{N}_m} w_j \cdot q_j (p) \cdot b_{jk} = 0
\end{equation*}

It is now useful to introduce a matrix $D$ of size $N \times N$ such that

\begin{equation*}
d_{ij} = \left\{
\begin{tabular}{ll}
1 & if $i,j$ produced by the same firm\\
0 & otherwise
\end{tabular}\right.
\end{equation*}

With $D$ we can write the previous first order condition as

\begin{equation*}
q_k(p) + \sum_{j \in \mathcal{N}_m} b_{jk} \cdot d_{jk} \cdot p_j - \sum_{j \in \mathcal{N}_m} b_{jk} \cdot d_{jk} \cdot v_j - 2 \cdot \sum_{j \in \mathcal{N}_m} b_{jk} \cdot d_{jk} \cdot w_j \cdot q_j (p) = 0
\end{equation*}

In this model each product is produced by exactly one firm. Therefore, we have one such first order condition for each product. Thus, we can write the system in a more compact way.

\begin{equation*}
a + B ^ t p + (D \circ B)(p - v - 2 \cdot w \cdot (a + B^t p)) = 0
\end{equation*}

Isolating for $p$, we obtain the optimum prices $p^*$

\begin{equation*}
p^* = [B^t + (D \circ B) - 2 (D \circ B)(w \circ B^t]^{-1} \cdot [-a + (D \circ B) v + 2 (D \circ B)(w \circ a)]
\end{equation*}

Then, the optimum quantites are defined by
\begin{equation*}
q^* = a + B^t \cdot p^*
\end{equation*}

We are now interested in computing consumer welfare $CS_j$ for each product. 
\begin{equation*}
CS_j = \int_{p_j^*}^{p_j(-\frac{a_j}{b_{jk})}} p_j d p_j = \frac{1}{2} \cdot [(p_j(-\frac{a_j}{b_{jk}}) - p_j^*) \cdot q_j^*]
\end{equation*}

\subsection{Quadratic Demands with Constant Marginal Costs}

Let $N$ be the number of products, and $\mathcal{N} = {1, 2, ..., N}.$
Let $q_k(p)$ be the demand of the k-th good. Because it is a quadratic demand, its curve is defined as 
\begin{equation*}
q_k(p) = a_k + \sum_{j=1}^N b_{kj}p_j + \sum_{j=1}^N e_{kj} p_j^2
\end{equation*}

Let $m$ be a firm tthat produces $\mathcal{N}_m \subseteq \mathcal{N}$. Then, its profit is defined by
\begin{equation*}
\pi_m(p) = \sum_{j \in \mathcal{N}_m} (p_j - c_j) \cdot q_j(p)
\end{equation*}

And its costs $C_j$ are characterized by 
\begin{equation*}
C_j(q_j) = f_j + vj \cdot q_j
\end{equation*}

Where $j \subseteq \mathcal{N}_m$. Thus, marginal costs are,
\begin{equation*}
c_j = \frac{\partial C_j(q_j)}{\partial q_j} = v_j
\end{equation*}

The firm, then, seeks to maximize its profit by setting their prices according to the following formula
\begin{equation*}
\max_{\{p_j ~ | ~ j \in \mathcal{N}_m\}} \pi_m(p)
\end{equation*}

If we derive with respect to $p_k$ where $k \in \mathcal{N}_m$, we obtain
\begin{eqnarray*}
\frac{\partial \pi_m(p)}{\partial p_k} &=& \frac{\partial}{\partial p_k} \sum_{j \in \mathcal{N}_m} (p_j - c_j) \cdot q_j(p)\\
&=& \sum_{j \in \mathcal{N}_m} \frac{\partial}{\partial	p_k} [(p_j - c_j) \cdot q_j(p)]\\
&=& \sum_{j \in \mathcal{N}_m]} \frac{\partial (p_j - c_j)}{\partial p_k} \cdot q_j(p) + \sum_{j \in \mathcal{N}_m} \frac{\partial q_j(p)}{\partial p_k} \cdot (p_j - c_j)\\
&=& q_j(p) + \sum_{j \in \mathcal{N}_m} (b_{jk} + 2 \cdot e_{jk} \cdot p_k) \cdot (p_j - c_j)\\
\end{eqnarray*}

From here we obtain the first order condition
\begin{equation*}
q_j(p) + \sum_{j \in \mathcal{N}_m} (b_{jk} + 2 \cdot e_{jk} \cdot p_k) \cdot (p_j - c_j) = 0
\end{equation*}

It is useful, now, to introduce a matrix $D$ of size $N \times N$ such that
\begin{equation*}
d_{ij} = \left\{
\begin{tabular}{ll}
1 & if $i,j$ produced by the same firm\\
0 & otherwise
\end{tabular}\right.
\end{equation*}

With $D$ we can write the previous first order condition as

\begin{equation*}
q_j(p) + \sum_{j \in \mathcal{N}_m} (d_{jk} \cdot b_{jk} + 2 \cdot d_{jk} \cdot e_{jk} \cdot p_k) \cdot (p_j - c_j) = 0
\end{equation*}

In this model, each product is produced by exactly one firm. Therefore, we have one such first order condition for each product. Thus, we can write the system in a more compact way.
\begin{equation*}
a +(B \circ D) p + (E \circ D) (p \circ p) + ((B \circ D) + 2 (E \circ D) p) (p - c) = 0 
\end{equation*}

Because this equation cannot be isolated for $p$, we obtain the optimum prices through the Newton-Raphson Method.

\begin{equation*}
P_{n+1} = P_n - J^{-1} \cdot F(P_n)
\end{equation*}

Where $J^{-1}$ is the inverse of the Jacobian matrix, and $F_j(P_n) = \partial \pi_m / \partial p_k$. The Jacobian matrix is constructed as follows

\begin{equation}
J_{ij} = \frac{\partial F_i}{\partial p_j} = \frac{\partial q_i(p)}{\partial p_j} + \sum_{j \in \mathcal{N}_m} \frac{\partial^2 q_i(p)}{\partial p_i \partial p_j} + \sum_{j \in \mathcal{N}_m} \frac{\partial q_i(p)}{\partial p_j}
\end{equation}

Then, the optimum quantities are defined by 
\begin{equation*}
q^* = a + B^t p^* + E^t p^{*2}
\end{equation*}

We are now interested in computing consumer welfare $CS_j$ for each product.
\begin{equation*}
CS_j = \int_{p_j^*}^{p_j(a + b_{jk} p_j + e_{jk} p_j^2 = 0)} p_j dp_j
\end{equation*}

To find $a + b_{jk} p_j + e_{jk} p_j^2 = 0$ it is necessary to use the Newton-Raphson method again. 

\subsection{Quadratic Demands with Linear Marginal Costs}

Let $N$ be the number of products, and $\mathcal{N} = 1, 2, ..., N$. Let $q_k(p)$ be the demand of the k-th good. Because it is a quadratic demand, its curve is defined as
\begin{equation*}
q_k(p) = a_k + \sum_{j = 1}^N b_{kj} p_j + \sum_{j = 1}^N e_{kj}p_j^2
\end{equation*}

Let $m$ be a firm that produces $\mathcal{N}_m \subseteq \mathcal{N}$. Then, its profit is defined by
\begin{equation*}
\pi_m(p) = \sum_{j \in \mathcal{N}_m} (p_j - c_j) \cdot q_j(p)
\end{equation*}

And its costs $C_j$ are characterized by
\begin{equation*}
C_j(q_j) = f_j + v_j \cdot q_j + w_j \cdot q_j^2
\end{equation*}

Where $j \subseteq \mathcal{N}_m$. Thus, marginal costs are,
\begin{equation*}
c_j = \frac{\partial C_j(q_j)}{\partial q_j}= v_j + w_j \cdot q_j
\end{equation*}

The firm, then, seeks to maximize its profit by setting their prices according to the following formula
\begin{equation*}
\max_{\{p_j ~ | ~ j \in \mathcal{N}_m\}} \pi_m(p)
\end{equation*}

If we derive with respect to $p_k$ where $k \in \mathcal{N}_m$, we obtain
\begin{eqnarray*}
\frac{\partial \pi_m(p)}{\partial p_k} &=& \sum_{j \in \mathcal{N}_m} \frac{\partial}{\partial p_k} (p_j - v_j - w_j \cdot q_j) q_j\\
&=& q_k - \sum_{j \in \mathcal{N}_m} \frac{\partial q_j}{\partial p_k} w_j \cdot q_j + \sum_{j \in \mathcal{N}_m} (p_j - v_j - w_j \cdot q_j) \frac{\partial q_j}{\partial p_k}\\
&=& q_k - \sum_{j \in \mathcal{N}_m} (b_{jk} + 2e{jk} \cdot p_j) w_j \cdot q_j + \sum_{j \in \mathcal{N}_m} (p_j - v_j - w_j \cdot q_j)(b_{jk} + 2e_{jk} \cdot p_k
\end{eqnarray*}

From here we obtain the first order condition
\begin{equation*}
q_k - \sum_{j \in \mathcal{N}_m} (b_{jk} + 2e{jk} \cdot p_j) w_j \cdot q_j + \sum_{j \in \mathcal{N}_m} (p_j - v_j - w_j \cdot q_j)(b_{jk} + 2e_{jk} \cdot p_k = 0
\end{equation*}

It is useful, now, to introduce a matrix $D$ of size $N \times N$ such that 
\begin{equation*}
d_{ij} = \left\{
\begin{tabular}{ll}
1 & if $i,j$ produced by the same firm\\
0 & otherwise
\end{tabular}\right.
\end{equation*}

With $D$ we can write the previous first order condition as
\begin{equation*}
q_k - \sum_{j \in \mathcal{N}_m} (b_{jk} \cdot d_{jk} + 2e{jk} \cdot d{jk} \cdot p_j) w_j \cdot q_j + \sum_{j \in \mathcal{N}_m} (p_j - v_j - w_j \cdot q_j)(b_{jk} \cdot d_{jk} + 2e_{jk} \cdot d_{jk} \cdot p_k = 0
\end{equation*}

In this model, each product is produced by exactly one firm. Therefore, we have one such first order condition for each product. Thus we can write the system in a more compact way.
\begin{eqnarray*}
a + B^t p + E^t (p \circ p) - ((B \circ D) &+& 2(E \circ D)p) (a + B^t p + E^t (p \circ p) w + \\
(p - v - w(a + B^t p &+& E^t (p \circ p)))((B \circ D) + 2 (E \circ D) p) = 0
\end{eqnarray*}

Because this equation cannot be isolated for $p$, we obtain the optimum prices through the Newton-Raphson Method.
\begin{equation*}
P_{n+1} = P_n - J^{-1} \cdot F(P_n)
\end{equation*}

Where $J^{-1}$ is the inverse of the Jacobian matrix, and $F_j(P_n) = \partial \pi_m / \partial p_k$. The Jacobian matrix is constructed as follows

\begin{eqnarray*}
J_{ij} = \frac{\partial F_i}{\partial p_j} = \frac{\partial q_i}{\partial p_j} &-& \sum_{k \in \mathcal{N}_m} \frac{\partial^2 q_k}{\partial p_i \partial p_j} w_k \cdot q_k \cdot d_{ik}\\
 - 2\sum_{k \in \mathcal{N}_m} w_k \cdot d_{ik} \frac{\partial q_k}{\partial p_i} \frac{\partial q_k}{\partial p_j} + \frac{\partial q_j}{\partial p_i} d_{ik} &+& \sum_{k \in \mathcal{N}_m} (p_k - v_k - w_k \cdot q_k) \frac{\partial^2 q_k}{\partial p_i \partial p_j}d_{ik}
\end{eqnarray*}

Then, the optimum quantities are defined by 
\begin{equation*}
q^* = a + B^t p^* + E^t p^{*2}
\end{equation*}

We are now interested in computing consumer welfare $CS_j$ for each product.
\begin{equation*}
CS_j = \int_{p_j^*}^{p_j(a + b_{jk} p_j + e_{jk} p_j^2 = 0)} p_j dp_j
\end{equation*}

To find $a + b_{jk} p_j + e_{jk} p_j^2 = 0$ it is necessary to use the Newton-Raphson method again. 


An example of reference is \cite{abc}.

\newpage
\begin{thebibliography}{9}

\bibitem{abc}
ABC: A System for Sequential Synthesis and Verification. Berkeley Logic Synthesis and Verification Group.
 
\bibitem{dc}
Design Compiler, from Synopsys. \url{http://www.synopsys.com/Tools/Implementation/RTLSynthesis/DesignCompiler/Pages/default.aspx}

\bibitem{cudd}
CUDD: CU Decision Diagram Package. Fabio Somenzi, Colorado University. \url{http://vlsi.colorado.edu/~fabio/CUDD/}.

\bibitem{elmore48}
W. C. Elmore. The Transient Response of Damped Linear Networks with Particular Regard to Wideband Amplifiers. Journal of Applied Physics, Vol. 19, 1948.

\bibitem{nangate_library}
NanGate Open Cell Library. \url{http://www.nangate.com/?page_id=22}.
\end{thebibliography}

\newpage
\listoffigures

\newpage
\listoftables

%
\end{document}
%--------------------------------------------------------------%
%--------------------------------------------------------------%
